\documentclass[10pt,a4paper]{amsart}
\usepackage[margin=19mm]{geometry}
\usepackage{amsmath,amssymb,mathtools}
\usepackage[scr=rsfs,cal=euler]{mathalfa}
\usepackage{xfrac}
\usepackage[unicode,hidelinks,bookmarks=false]{hyperref}
\newcommand{\ie}{i.e.\ }
\newcommand{\cf}{cf.\ }
\newcommand{\resp}{respectively\ }
\newcommand{\Subsection}{\subsection}
\providecommand{\nobreakdash}{\nobreak\hbox{-}}
\setlength{\emergencystretch}{2em}
\multlinegap0pt
\usepackage{amssymb}
\usepackage[shortlabels]{enumitem}
\usepackage{tikz}
\usepackage{float}
\usepackage{mathtools}
\newtheorem{theorem}{Theorem}
\newcommand{\SO}{\mathrm{SO}}
\newcommand{\TT}{\mathcal{T}}
\newcommand{\irr}{\mathrm{irr}}
\title{Theta-Relations Among Degree-Based Tree Indices}
\author{Duaa Abdullah, Jasem Hamoud}
\date{Source: arXiv:2602.21318v2 (CC BY 4.0)}
\begin{document}
\maketitle

\noindent\emph{Source and licence.} Theta-Relations Among Degree-Based Tree Indices. Version arXiv:2602.21318v2 (CC BY 4.0), distributed by arXiv under the Creative Commons Attribution 4.0 International licence, http://creativecommons.org/licenses/by/4.0/, as recorded in the arXiv licence metadata for this exact version and shown on the abstract page https://arxiv.org/abs/2602.21318v2. Full licence notice: \texttt{licenses/arxiv-2602.21318-CC-BY-4.0.txt}.\par\smallskip

\noindent\textit{Editorial scope.} Counted unit: the article's theorem ThmSomborn3 (source0.tex lines 549--558). The article's complete original proof of it is delivered with it (source0.tex lines 559--616, one \texttt{proof} environment of 58 lines). No mathematics is written, repaired or re-derived: the statement and the proof are the article's own lines, in the article's own notation, and no retained line is substituted or re-flowed.  Retained uncounted as the source-local context the proof needs: lines 384--386; lines 490--496.  Nothing was omitted from any retained span, so the package carries no omission record.  No retained line contains a citation, so the wrapper carries no bibliography and no bibliography entry.  No retained line contains a figure environment, an image-inclusion command or drawing code, so no image is delivered and the package declares no asset dependency.  Editorial formatting only, in addition: an AMS \texttt{amsart} preamble that restates the article's own declarations and macros used by the retained lines, namely the environments \texttt{theorem} on separate counters, together with the standard packages those lines need.  The retained environments print in this document as Theorem 1 in order of appearance, whereas the article prints the same items with the numbering of its own full document.  The complete native source of the article as distributed in the arXiv e-print for this exact version is bundled unchanged as \texttt{sources/195262/source0.tex}.
\begin{equation}~\label{eqqq4lemSomborIn2}
\sigma(\TT)=\sum_{v\in V(\TT)}d_\TT(v)^2-\frac{4(n-1)^2}{n}.
\end{equation}

\begin{theorem}~\label{ThmSomborn2}
Let $\TT\in \mathcal{T}_{n,\Delta}$ be a tree of order $n$.
Then
\begin{equation}~\label{eqq1ThmSomborn2}
\frac{1}{\sqrt{2}}\Bigl(\sigma(\TT)+\frac{4m^2}{n}\Bigr) \leqslant \SO(\TT)\leqslant\Bigl(\sigma(\TT)+\frac{4m^2}{n}\Bigr).
\end{equation}
\end{theorem}

\begin{theorem}~\label{ThmSomborn3}
Let $\TT$ be a tree of order $n$. Then
\begin{equation}~\label{eqq1ThmSomborn3}
\irr(\TT)\leqslant \SO(\TT)\leqslant\sqrt{2}\,\irr(\TT)+\sum_{v \in V(T)} d(v)^2.
\end{equation}
Moreover,
\[
\SO(\TT)=\Theta\!\left(\irr(\TT)+\sum_{v\in V(\TT)} d(v)^2\right).
\]
\end{theorem}

\begin{proof}
Assume $\TT\in \mathcal{T}_{n,\Delta}$ be a tree of order $n$. According to Theorem~\ref{ThmSomborn2} the relationship between Sigma index and Sombor index satisfy $\sigma(\TT)\leqslant \SO(\TT)$. Since $\sqrt{\sigma(\TT)}\leqslant \irr(\TT)\leqslant \sqrt{m\,\sigma(\TT)}$ implies that $\sqrt{\SO(\TT)}\leqslant \irr(\TT)\leqslant \sqrt{2n\,\SO(\TT)}$. Consider an arbitrary edge $uv \in E(\TT)$. The corresponding term in $\SO(\TT)$ is $\sqrt{d(u)^2+d(v)^2}$. For any non-negative real numbers $x$ and $y$, the inequality $|x-y|\leqslant \sqrt{x^2+y^2}$ holds, since $(x-y)^2=x^2-2xy+y^2\leqslant x^2+y^2$ follows from $-2xy\leqslant 0$ it yields $|d(u)-d(v)|\leqslant \sqrt{d(u)^2+d(v)^2}$. Thus, the lower bound of Sombor index among~\eqref{eqq1ThmSomborn3} holds.

To prove the upper bound, we find in~\eqref{eqqq4lemSomborIn2} the following relationship
\[
\sum_{v\in V(\TT)}d_\TT(v)^2=\sigma(\TT)+\frac{4(n-1)^2}{n}.
\]
Thus, for any edges $uv\in E(\TT)$ it follows that $|d(u)-d(v)|\leqslant (d(u)+d(v))\sqrt{d(u)^2+d(v)^2}$  and the Albertson index satisfy
\begin{equation}~\label{eqq2ThmSomborn3}
\irr(\TT)\geqslant \sqrt{\sigma(\TT)+\frac{4(n-1)^2}{n}}.
\end{equation}
 Consider $\lambda_c=\max\{m\,\sigma(\TT),2n\,\SO(\TT)\}$. Then, from~\eqref{eqq2ThmSomborn3} we observe that $\irr(\TT)\leqslant \lambda_c+\frac{4(n-1)^2}{n}$. For any edge $uv \in E(\TT)$, assume $t_1=(d(u)-d(v))^2+(d(u)+d(v))^2$ the following identity holds $d(u)^2+d(v)^2=t_1/2$. It follows that $\sqrt{d(u)^2+d(v)^2}=\sqrt{t_1/2}$  yields
 \begin{align*}
 \sqrt{d(u)^2+d(v)^2} &\leqslant \frac{1}{\sqrt{2}}(|d(u)-d(v)|+d(u)+d(v))\\
 &\leqslant \frac{\lambda_c}{\sqrt{2}}(2\,|d(u)-d(v)|)\\
 &\leqslant \sqrt{2}\,|d(u)-d(v)|.
 \end{align*}
By induction, the upper bound on the Sombor index stated in \eqref{eqq1ThmSomborn3} holds for $n=2$. For $n=3$, we have $\mathcal{P}_n \cong \mathcal{S}_n$, and thus the same upper bound holds. Assume the inequality holds for all trees of order $k<n$, where $n \geqslant 3$. Let $w$ be a leaf adjacent to vertex $v$, and let $\TT'=\TT-w$, which is a tree of order $n-1$. Let $d=d_{\TT'}(v)$; then $d_\TT(v)=d+1$ and $d_\TT(w)=1$. The invariants transform as follows $\irr(\TT)=\irr(\TT')+|d_\TT(v)-d_\TT(w)|$, $\irr(\TT')=\irr(\TT)-d$, $\SO(\TT)=\SO(\TT')+\sqrt{d_T(v)^2 + d_T(w)^2}$, $\SO(\TT')=\SO(\TT)-\sqrt{(d+1)^2+1}$, $M(\TT)=M(\TT')+d_\TT(v)^2+d_\TT(w)^2-d_{\TT'}(v)^2$ and $M(\TT')=M(\TT)-2d-2$. Therefore, by considering the equation~\eqref{eqq2ThmSomborn3}, it follows that $\SO(\TT')\leqslant \irr(\TT')+M(\TT')$. Thus,  
\begin{equation}~\label{eqq3ThmSomborn3}
\SO(\TT')+\sqrt{(d+1)^2+1}\leqslant \sqrt{2}\, \irr(\TT')+M(\TT')+ \sqrt{(d+1)^2+1}.
\end{equation}
Hence, we have 
\begin{equation}~\label{eqq4ThmSomborn3}
\sqrt{2}\,\irr(\TT)+M(\TT)=\sqrt{2}\,(\irr(\TT')+d)+M(\TT')+2d+2.
\end{equation}
Since $d \geqslant 1$, from~\eqref{eqq3ThmSomborn3} and \eqref{eqq4ThmSomborn3} implies that
\[
\sqrt{(d+1)^2+1}=(d+1)\sqrt{1+\frac{1}{(d+1)^2}}.
\]
By the concavity of the square root function, its graph lies below the tangent line, which implies that for all $y>0$,$\sqrt{1+y}<1+\frac{y}{2}$. Thus, by considering $y=\frac{1}{(d+1)^2}$ yields $\sqrt{1+\frac{1}{(d+1)^2}}<1+\frac{1}{2(d+1)^2}$. Since $d+1>0$ gives
\begin{align*}
\sqrt{(d+1)^2+1}&= (d+1)\sqrt{1+\frac{1}{(d+1)^2}}\\
&<(d+1) \left(1+\frac{1}{2(d+1)^2}\right)\\
&=d+1+\frac{1}{2(d+1)}.
\end{align*}
Since $d \geqslant 1$, we have $d+1+\frac{1}{2(d+1)}\leqslant d+1+\frac{1}{4}$. 

\vspace{0.3 cm}
\noindent\textbf{The Asymptotic Relation:} $\SO(\TT)=\Theta\!\bigl(\irr(\TT)+M(\TT)\bigr)$. Assume there exist positive constants $a_1$ and $a_2$, independent of $\TT$ such that
\begin{equation}~\label{eqq5ThmSomborn3}
a_1 \bigl(\irr(\TT)+M(\TT)\bigr) \leqslant \SO(\TT) \leqslant a_2 \bigl(\irr(\TT)+M(\TT)\bigr)
\end{equation}
From a previously~\eqref{eqq3ThmSomborn3} and \eqref{eqq4ThmSomborn3} established inequality. 
Hence $a_2=\sqrt{2}$ is sufficient. We first note the trivial inequality $\SO(\TT) \geqslant \irr(\TT)$.
Then, we derive a lower bound in terms of $M(\TT)$. For every edge $uv \in E(\TT)$, $\sqrt{d(u)^2+d(v)^2}\geqslant \frac{d(u)+d(v)}{\sqrt{2}}$. Then, 
\begin{equation}~\label{eqq6ThmSomborn3}
\SO(\TT)\geqslant \frac{1}{\sqrt{2}}\,\sum_{uv\in E(\TT)}(d_\TT(u)+d_\TT(v)).
\end{equation}
Therefore, from~\eqref{eqq5ThmSomborn3} and \eqref{eqq6ThmSomborn3} we obtain 
$\SO(\TT)\geqslant \sqrt{2}\,M(\TT)$. 
If $\irr(\TT)$ dominates (for example, in stars $\mathcal{S}_n$), then $\SO(\TT)$ is asymptotically close to $\sqrt{2}\,\bigl(\irr(\TT)+M(\TT)\bigr)$. If $M(\TT)$ dominates (for example, in paths $\mathcal{P}_n$), then $\SO(\TT) \geqslant\sqrt{2}\, M(\TT)>\frac{1}{2} \bigl(\irr(\TT)+M(\TT)\bigr)$. Then, 
\begin{equation}~\label{eqq7ThmSomborn3}
\SO(\TT)=(n-1)^2\,\sqrt{1+\frac{1}{(n-1)^2}}
\end{equation}
Thus, from~\eqref{eqq7ThmSomborn3}, we have  $\SO(\mathcal{S}_n)\geqslant \irr(\TT)+M(\TT)$.
For the path $\mathcal{P}_n$, $\SO(\mathcal{P}_n)\geqslant 2\sqrt{2}\,n$.
Hence,  $\irr(\TT)+M(\TT)\geqslant 4n$. Thus, $\SO(\mathcal{S}_n)\geqslant \left(\irr(\TT)+M(\TT)\right)/2$. As desire.
\end{proof}

\end{document}
